\section{课程机制：作业、AI policy 与学习方式}

CS336 是 5-unit 课程，作业强度非常高。讲义引用 Spring 2024 评价说，第一份作业接近 CS224N 五份作业加 final project 的工作量。这不是为了堆工作，而是为了让学生真正实现和 benchmark。

\begin{longtable}{p{0.20\textwidth}p{0.55\textwidth}p{0.15\textwidth}}
\toprule
\textbf{作业} & \textbf{内容} & \textbf{单元} \\
\midrule
Assignment 1 & BPE tokenizer、Transformer、cross-entropy、AdamW、training loop、TinyStories/OpenWebText leaderboard & Basics \\
Assignment 2 & fused RMSNorm kernel、distributed data parallel、optimizer state sharding、benchmark/profile & Systems \\
Assignment 3 & 小规模训练 API、fit scaling laws、预测大规模 loss/超参 & Scaling \\
Assignment 4 & Common Crawl 到文本、filtering、dedup、data mixing、perplexity leaderboard & Data \\
Assignment 5 & DPO、GRPO、偏好优化和 RL systems & Alignment \\
\bottomrule
\end{longtable}

\begin{warningbox}{AI policy 的真实含义}
Coding agents 可能能完成作业，但如果学生只拿答案，就绕过了课程目标。AI 可以用来解释、提问、debug，但不能替代理解 mechanics。讲义生成也一样：不能只给结论，要展示机制和推理。
\end{warningbox}

\subsection{本章小结}

课程机制服务于学习目标：通过实现、测试、benchmark 和资源核算建立理解。作业不是普通 coding exercise，而是逐步搭建语言模型训练系统。

